\section{Conclusion and Future Work}
\label{sec:conclusion}

In this work, we introduced \textbf{PAC-Kinetics}, the first mathematically rigorous framework for stateful, sequential model ensembling. By integrating continuous-time non-equilibrium chemical kinetics with PAC-Bayesian generalization theory for stationary $\beta$-mixing stochastic processes, we resolved the standard accuracy-stability trade-off that plagues test-time adapter serving. We provided contractive and Lyapunov stability proofs demonstrating that our dynamical router is globally asymptotically stable and input-to-state stable (ISS). 

Through comprehensive experiments in our Coordinates Sandbox, PAC-Kinetics achieved a massive reduction in routing jitter—over \textbf{11.2$\times$} on orthogonal streams and up to \textbf{16.0$\times$} on overlapping streams compared to stateless SABLE—while matching or exceeding Oracle accuracy. Furthermore, unlike heuristic stateful approaches that collapse under heterogeneous stream switches, our PAC-Bayesian optimized parameters remain robustly stable.

Our work bridges the gap between machine learning theory, dynamical systems, and real-time edge serving.

\subsection{Limitations and Path to Physical Verification}
While our work establishes a rigorous mathematical foundation, we acknowledge that our empirical validation is restricted to the simulated Coordinates Sandbox (ICS). In the ICS, representation trajectories are simulated inside a closed-form vector space, and task classification is evaluated via a soft continuous accuracy metric. This simplification is highly valuable for theoretical research: it is the only way to precisely control and verify mixing coefficients $\beta(j)$ and coordinate noise scales, which is critical for confirming that our PAC-Bayesian bound works exactly as proved.

However, deploying PAC-Kinetics on real-world networks is our ultimate goal. In physical deployments (e.g., serving multi-task Vision Transformers or Large Language Models using LoRA adapters on standard datasets like GLUE or Decathlon), several physical realities must be managed. First, intermediate representation scales may fluctuate across layers, requiring adaptive layer-wise normalization. Second, the wall-clock latency of tracking the routing state $s_t$ must be negligible compared to the backbone's execution latency. In our preliminary tests, maintaining a 4-dimensional state $s_t$ adds less than 50 microseconds of overhead, which is completely negligible compared to the millisecond-scale latency of neural forward passes. 

Future directions include deriving generalization bounds for more general piecewise-stationary or non-stationary request streams, and extending this continuous-time stateful ensembling paradigm to physical, large-scale generative language and multi-modal models. 

\paragraph{Open Challenge: Scaling the Physical Expert Fleet} Explicitly evaluating physical networks under larger expert fleet sizes (such as $K=4$ or $K=8$) on complex datasets is a highly recommended direction for subsequent systems research. Crucially, while we have conducted systems-level latency profiling on CPU for fleet sizes $K \le 8$, evaluating PAC-Kinetics on massive Transformer backbones (such as LLMs or ViTs with dozens of layers) under large fleet scales remains a critical open direction for future systems research.

\paragraph{Systems-Level GPU Serving and Multi-Batching Overhead} While our current systems profiling is conducted on a single CPU core (where PAC-Kinetics exhibits flat $O(1)$ execution latency of $\approx 10.4$ microseconds per sample), deploying stateful routers on high-throughput GPU serving infrastructure introduces unique systems-level considerations. In active multi-batching serving frameworks (such as S-LoRA or Punica), queries are processed in large concurrent batches. To support this on GPU hardware, the stateful linear recurrence $s_t = \mathbf{A} s_{t-1} + W \mathbf{e}_t$ must be implemented as a custom CUDA kernel or vectorized PyTorch tensor operation. Since the recurrence is linear and element-wise across the batch, it can be executed in parallel with virtually zero CUDA scheduling overhead. Specifically, tracking and updating the concentration state $s_t$ across a batch of size 128 under $K=16$ experts requires less than 2,048 floating-point operations (FLOPs), which can be completed in less than a microsecond on modern tensor cores. This is several orders of magnitude smaller than the millisecond-scale execution latency of the underlying attention layers, confirming that PAC-Kinetics is highly compatible with high-throughput GPU serving systems.

\subsection{Online Adaptability and Real-Time Tuning}
An exciting avenue of future work is extending PAC-Kinetics to dynamically adjust its kinetics parameters online without requiring explicit re-calibration splits. Specifically, the system could estimate the real-time temporal correlation of incoming queries—for example, by computing a running autocorrelation of the coordinates $\mathbf{e}_t$. When the autocorrelation is high (indicating a homogeneous task block), the router could dynamically increase the state retention coefficients $a_k$ to maximize noise filtering and smoothness. Conversely, when the autocorrelation falls (indicating rapid task switching or independent queries), the system could reduce $a_k$ to suppress history and avoid ``inertial drag,'' allowing the router to adapt instantly to switches. This online adaptive kinetics loop would automatically traverse the accuracy-stability Pareto frontier in real-time, providing optimal performance across arbitrary, dynamically shifting workloads.

\paragraph{Bridging Theoretical Mixing Bounds to Online Autocorrelation} In production edge serving environments, the true joint distribution of incoming queries is typically unknown, which makes the mixing coefficient $\beta(b)$ and the corresponding failure probability term $1 - \delta - 2(a/2-1)\beta(b)$ in Theorem 3.1 practically unverifiable. To bridge this gap between theory and practice, systems administrators can utilize our proposed online coordinate autocorrelation tracking mechanism as a qualitative proxy for the mixing dynamics. Specifically, by computing a rolling autocorrelation of the incoming PCA coordinates $\mathbf{e}_t$, practitioners can dynamically monitor the temporal dependency structure of the stream. A high rolling autocorrelation indicates low mixing (high temporal dependency, corresponding to slow-decaying $\beta(j)$ coefficients), signaling that the router should operate with high state retention. Conversely, a drop in autocorrelation indicates rapid mixing (corresponding to fast-decaying $\beta(j)$ coefficients), signaling that the failure probability bound is mathematically tight and the system can safely operate with lower retention to avoid inertial drag. This provides a concrete systems-level bridge from our abstract learning-theoretic probability bounds to active, real-time workload monitoring.

\subsection{Learnable Parameter Uncertainty under PAC-Bayes}
In our current formulation, we fixed the posterior covariance matrix to match the prior default $\sigma_0^2 I$, which elegantly simplifies the Gaussian KL complexity penalty to classical $L_2$ parameter distance (making it equivalent to standard weight decay centered at SABLE defaults). While this simplification is highly valuable as it bridges our learning-theoretic bound to classical regularized optimization practices, it discards a major benefit of PAC-Bayesian theory: the ability to capture parameter uncertainty.

An exciting avenue of future work is to optimize a diagonal vector of learnable posterior variances $\boldsymbol{\sigma}^2 = [\sigma_1^2, \dots, \sigma_M^2]^T$ along with the parameter mean vector $\Theta$. Optimizing the posterior variance allows the gradient engine to automatically learn which parameters are highly sensitive to coordinate noise and restrict their updates, while relaxing constraints on highly robust dimensions. This learnable parameter uncertainty could yield significantly tighter generalization bounds and even greater serving stability, providing a theoretically unified mechanism for managing noise sensitivity in stateful routers.
